\scanpage{24}
\mspage{7}
of $\dot X\subset\widetilde X$ is determined by knowledge of
\[
\widetilde X'_1\cap\dot X=\dot X'_1.
\]

\uline{Corollary} The \inserted{closed} subsets $X'$ of $X$ compatible with
the germ of conical structure, are in
$1$-$1$ correspondence with the pairs $(\widetilde X'\subset\widetilde X,Y'\subset Y)$
of closed subsets of $\widetilde X$ and of $Y$, such that

\unsure{a)|1°)} $\widetilde X'$ is compatible with the \inserted{germ of} cylindrical structure
\inserted{of $\widetilde X,\dot X$}

\unsure{b)|2°)} $f:\dot X\to Y$ maps \deleted{$\unsure{\widetilde X'|\widetilde X}$}
$\dot X'\overset{\text{definition}}{=}\widetilde X'\cap\dot X$ into $Y'$.

We shall have
% P07: the upper brace begins beneath the right half of Im's m;
% it is represented as attaching to the whole image expression.
\[
X'=\underbrace{\operatorname{Im}(\widetilde X'\ \text{under}\ \widetilde X\to X)}
_{\overline{X'\setminus\underbrace{Y\cap X'}_{Y'}}}\cup Y'
\]
For $X'$ to be \deleted{\uline{strictly\unsure{.}} compatible with}
\inserted{\uline{strictly} conical},
% P09: the cancelled raised word attaches after the old article, before structure.
\deleted{\unsure{the|the} \inserted{\unsure{purely|germ of}} conical structure}
it is necessary and sufficient that $Y'=\operatorname{Im}(\dot X'\to Y)$

Let now $X_0\subset X$ be a
bounding subset of $X$, and $\mathcal E$ a \inserted{germ of} cylindrical structure
of $X_0$ in $X$. We say that $\mathcal E$ is compatible
with the germ of conical structure $\mathcal T$
of $X$ around $Y$, if $\exists$ representatives
$E\inserted{\simeq X_0\times\mathbf I}$ cylindrical neighborhood of $X_0$, \unsure{\&|and} $T$ conical neighborhood
of $Y$ (with \inserted{\unsure{its|the}} conical structure), such that

\margin{instead of ``compatible'' it is more telling to
say ``\uline{transverse}''}

a) $T\cap\unsure{\overline E|E}$, $\dot T\cap E$, and $Y\cap\unsure{\overline E|E}$
\deleted{\unsure{compatible}} \inserted{are} cylindrical subsets of
\deleted{$\unsure{\overline E|E}\cap\unsure{\overline T|T}$ is stable \unsure{under|under} the two types of \unsure{homotheties}
$\unsure{\overline x|x}\mapsto\unsure{{\lambda x}_{X_0}|\lambda_{X_0}x}$ and
$t\unsure{\longrightarrow|\mapsto}\unsure{\lambda_Y\times t|\lambda\times_Y t}$
\unsure{which|that} \unsure{act} \unsure{on either side}}
% P18: the second map's Y is lowered between lambda and the crossing sign;
% attachment to lambda is preferred but not settled. The final t is crossed.
\deleted{$E$ \unsure{i.e.}}

\margin{hence stable under $E$-homotheties for $\lambda\in[0,1]$}

% P20: the entire diagram is cancelled by nine long diagonal strokes.
% The leftmost source stroke also crosses the separately transcribed E i.e.
% Normal diagram spacing is used; no extra arrowheads are supplied to the cups.
\diagram{
\begin{tikzcd}[ampersand replacement=\&,
 column sep=large,row sep=large,
 execute at end picture={
  \draw[line width=.6pt] ([xshift=-2.4cm,yshift=.5cm]pA.north west) -- ([xshift=-1.1cm,yshift=-.4cm]pF.south west);
  \draw[line width=.6pt] ([xshift=-1.6cm,yshift=.8cm]pA.north west) -- ([xshift=-.2cm,yshift=-.3cm]pF.south west);
  \draw[line width=.6pt] ([xshift=-.8cm,yshift=.8cm]pA.north west) -- ([xshift=.6cm,yshift=-.4cm]pF.south west);
  \draw[line width=.6pt] ([xshift=-.3cm,yshift=.8cm]pA.north west) -- ([xshift=1.0cm,yshift=-.1cm]pF.south west);
  \draw[line width=.6pt] ([xshift=.2cm,yshift=.7cm]pA.north east) -- ([xshift=1.2cm,yshift=-.5cm]pF.south east);
  \draw[line width=.4pt] ([xshift=-.6cm,yshift=.8cm]pB.north west) -- ([xshift=.6cm,yshift=-.1cm]pG.south west);
  \draw[line width=.6pt] ([xshift=.1cm,yshift=.6cm]pB.north) -- ([xshift=1.5cm,yshift=-.5cm]pG.south);
  \draw[line width=.6pt] ([xshift=-.6cm,yshift=1.6cm]pE.north west) -- ([xshift=.5cm,yshift=-.2cm]pH.south west);
  \draw[line width=.6pt] ([xshift=.2cm,yshift=1.6cm]pE.north) -- ([xshift=1.4cm,yshift=-.4cm]pH.south);
 }
]
|[alias=pA]| E \arrow[r,phantom,"\simeq"] \&
|[alias=pB]| X_0\times I \\
|[alias=pC]| \dot T\cap\unsure{E|\overline E}
 \arrow[u,phantom,"\cup" description] \arrow[r,phantom,"\simeq"] \&
|[alias=pD]| T_0\times I \arrow[u,phantom,"\cup" description] \&
|[alias=pE]| T_0=X_0\cap T \\
|[alias=pF]| Y\cap\unsure{E|\overline E} \arrow[r,phantom,"\simeq"] \&
|[alias=pG]| Y_0\times I \arrow[u,phantom,"\cup" description] \&
|[alias=pH]| Y_0=\dot X_0\cap Y
\end{tikzcd}}
% P20: a separately crossed small group lies to the left of the lower diagram.
\deleted{$\unsure{E_0I|E\cap T}$}
\margin{23}
